\begin{notechapter}
\chapter{The Singular Value Decomposition}\label{ch:svd}
\chaptersubtitle{Rotate, stretch, rotate}

\begin{chapterintro}
Every matrix, square or not, acts in the same simple way: it rotates (or reflects)
its input, stretches it along perpendicular axes, and rotates the result into the
output space. The singular value decomposition (SVD) makes this explicit, and it is
the computational backbone of PCA, least squares, low-rank approximation and
recommender systems. You will learn
\begin{itemize}[nosep]
  \item the statement of the SVD and a proof of its existence from the spectral
        theorem;
  \item its geometric meaning, and the reading of $A$ as a sum of rank-one maps;
  \item how to compute an SVD by hand on a $2\times2$ example;
  \item the best low-rank approximation (Eckart--Young);
  \item the Moore--Penrose pseudoinverse and the minimum-norm least-squares
        solution.
\end{itemize}
\prerequisites{\cref{ch:adjoint-operator}; eigenvalues and the spectral theorem for
symmetric matrices.}
\end{chapterintro}

\section{The decomposition}\label{svd:sec:statement}

\begin{theorem}[Singular value decomposition]\label{svd:thm:svd}
Every $A\in\R^{m\times n}$ of rank $r$ can be written
\[
  A = U\Sigma V^\top = \sum_{i=1}^r\sigma_i\,u_iv_i^\top,
\]
where $U = [u_1,\dots,u_m]\in\R^{m\times m}$ and $V = [v_1,\dots,v_n]\in\R^{n\times n}$ are
orthogonal, and $\Sigma\in\R^{m\times n}$ is zero except for its diagonal entries
$\Sigma_{ii} = \sigma_i$, with $\sigma_1\geq\sigma_2\geq\dots\geq\sigma_r>0$ and $\sigma_i = 0$
for $i>r$. The $\sigma_i$ are the \emph{singular values}, the $u_i$ the \emph{left} and the
$v_i$ the \emph{right singular vectors}. They satisfy
\[
  Av_i = \sigma_iu_i,\qquad A^\top u_i = \sigma_iv_i,\qquad
  A^\top Av_i = \sigma_i^2v_i,\qquad AA^\top u_i = \sigma_i^2u_i .
\]
\end{theorem}

\begin{proof}
$A^\top A$ is symmetric positive semi-definite ($x^\top A^\top Ax = \norm{Ax}^2\geq0$), so
by the spectral theorem it has an orthonormal basis of eigenvectors $v_1,\dots,v_n$ with
eigenvalues $\lambda_1\geq\dots\geq\lambda_n\geq0$. Since $\ker(A^\top A) = \ker(A)$
(\cref{adjoint-operator:prop:unique}), exactly $r$ eigenvalues are positive. Set
$\sigma_i = \sqrt{\lambda_i}$ and, for $i\leq r$, $u_i = Av_i/\sigma_i$. These are
orthonormal:
\[
  u_i^\top u_j = \frac{v_i^\top A^\top Av_j}{\sigma_i\sigma_j} = \frac{\lambda_j\,v_i^\top v_j}{\sigma_i\sigma_j} = \delta_{ij}.
\]
Complete $u_1,\dots,u_r$ into an orthonormal basis $u_1,\dots,u_m$ of $\R^m$. For $i>r$,
$\norm{Av_i}^2 = \lambda_i = 0$, so $Av_i = 0$. Hence $Av_i = \sigma_iu_i$ for every $i$
(with $\sigma_i = 0$ for $i>r$), that is $AV = U\Sigma$, and $A = U\Sigma V^\top$ because $V$
is orthogonal. The remaining relations follow by transposing and multiplying.
\end{proof}

The last $n-r$ columns of $V$ form an orthonormal basis of $\ker(A)$, and the last
$m-r$ columns of $U$ one of $\ker(A^\top) = \operatorname{Im}(A)^\perp$. The singular
values are unique; the singular vectors are unique up to sign when the $\sigma_i$ are
distinct. The \emph{thin} SVD keeps only the first $r$ columns: $A = U_r\Sigma_rV_r^\top$
with $U_r\in\R^{m\times r}$, $\Sigma_r = \diag(\sigma_1,\dots,\sigma_r)$, $V_r\in\R^{n\times r}$.

\section{Geometric meaning}\label{svd:sec:geometry}

Read $Ax = U\bigl(\Sigma(V^\top x)\bigr)$ from right to left:
\begin{enumerate}
  \item $V^\top x$ expresses $x$ in the orthonormal basis $v_1,\dots,v_n$: its entries are
        the coordinates $v_i^\top x = \norm x\cos\theta_i$ of $x$ along the input directions
        (a rotation or reflection of $\R^n$);
  \item $\Sigma$ multiplies coordinate $i$ by $\sigma_i$, drops the coordinates beyond $r$,
        and changes the dimension from $n$ to $m$;
  \item $U$ places the result along the output directions $u_1,\dots,u_m$ (a rotation or
        reflection of $\R^m$).
\end{enumerate}
Equivalently, by the rank-one sum,
\[
  Ax = \sum_{i=1}^r\sigma_i\,(v_i^\top x)\,u_i :
\]
each term \emph{projects} $x$ on $v_i$, \emph{scales} the coordinate by $\sigma_i$, and
\emph{embeds} it along $u_i$. In particular $A$ maps the unit sphere of $\R^n$ onto an
ellipsoid in $\operatorname{Im}(A)$ with semi-axes $\sigma_iu_i$
(\cref{svd:fig:ellipse}), and $\norm{A}_2 = \max_{\norm x=1}\norm{Ax} = \sigma_1$.

\begin{figure}[ht]
\centering
\begin{tikzpicture}[>=Stealth]
\begin{axis}[name=in, width=5.2cm, height=5.2cm, axis equal image, xmin=-1.6, xmax=1.6,
  ymin=-1.6, ymax=1.6, axis lines=middle, xtick=\empty, ytick=\empty, title={input $\R^2$}]
  \addplot[thick, bookaccent, domain=0:360, samples=90] ({cos(x)}, {sin(x)});
  \draw[->, very thick, bookexample] (axis cs:0,0) -- (axis cs:0.7071,0.7071) node[above right] {$v_1$};
  \draw[->, very thick, bookdef] (axis cs:0,0) -- (axis cs:0.7071,-0.7071) node[below right] {$v_2$};
\end{axis}
\begin{axis}[at={(in.east)}, anchor=west, xshift=2.2cm, width=7cm, height=7cm, axis equal image,
  xmin=-7, xmax=7, ymin=-7, ymax=7, axis lines=middle, xtick=\empty, ytick=\empty, title={output $\R^2$}]
  \addplot[thick, bookaccent, domain=0:360, samples=120] ({3*cos(x)}, {4*cos(x)+5*sin(x)});
  \draw[->, very thick, bookexample] (axis cs:0,0) -- (axis cs:2.1213,6.3640) node[right] {$\sigma_1u_1$};
  \draw[->, very thick, bookdef] (axis cs:0,0) -- (axis cs:2.1213,-0.7071) node[below right] {$\sigma_2u_2$};
\end{axis}
\draw[->, thick] ([xshift=4pt]in.east) -- node[above] {$A$} ([xshift=2.0cm]in.east);
\end{tikzpicture}
\caption{The matrix $A = \begin{psmallmatrix}3&0\\4&5\end{psmallmatrix}$ maps the unit circle to
an ellipse. The orthogonal input directions $v_1$, $v_2$ go to the orthogonal semi-axes
$\sigma_1u_1$, $\sigma_2u_2$ (\cref{svd:ex:2x2}).}
\label{svd:fig:ellipse}
\end{figure}

\begin{example}[An SVD by hand]\label{svd:ex:2x2}
Let $A = \begin{psmallmatrix}3&0\\4&5\end{psmallmatrix}$.
\begin{enumerate}
  \item $A^\top A = \begin{psmallmatrix}25&20\\20&25\end{psmallmatrix}$ has characteristic polynomial
        $(25-\lambda)^2 - 400$, so $\lambda_1 = 45$, $\lambda_2 = 5$, and
        $\sigma_1 = \sqrt{45} = 3\sqrt5\approx6.71$, $\sigma_2 = \sqrt5\approx2.24$.
  \item Eigenvectors: $v_1 = \frac1{\sqrt2}(1,1)$ and $v_2 = \frac1{\sqrt2}(1,-1)$.
  \item $u_1 = Av_1/\sigma_1 = \frac{1}{\sqrt2}(3,9)/(3\sqrt5) = \frac1{\sqrt{10}}(1,3)$ and
        $u_2 = Av_2/\sigma_2 = \frac1{\sqrt2}(3,-1)/\sqrt5 = \frac1{\sqrt{10}}(3,-1)$, which are indeed
        orthonormal.
\end{enumerate}
Hence
\[
  \begin{bmatrix}3&0\\4&5\end{bmatrix}
  = \underbrace{\frac1{\sqrt{10}}\begin{bmatrix}1&3\\3&-1\end{bmatrix}}_{U}
    \underbrace{\begin{bmatrix}3\sqrt5&0\\0&\sqrt5\end{bmatrix}}_{\Sigma}
    \underbrace{\frac1{\sqrt2}\begin{bmatrix}1&1\\1&-1\end{bmatrix}}_{V^\top}.
\]
Check: $\det A = 15 = \sigma_1\sigma_2$, and $\norm A_F^2 = 9+16+25 = 50 = \sigma_1^2+\sigma_2^2$.
\end{example}

\begin{pitfall}[Singular values are not eigenvalues]
The eigenvalues of this $A$ are $3$ and $5$ (it is triangular), not $6.71$ and $2.24$.
Eigenvectors need not be orthogonal and describe directions that $A$ maps to
themselves; singular vectors are always orthonormal and describe how $A$ maps one
orthonormal frame to another. The two coincide (up to sign) only for symmetric positive
semi-definite matrices.
\end{pitfall}

\section{Best low-rank approximation}\label{svd:sec:eckart-young}

Truncating the rank-one sum after $k$ terms gives
$A_k = \sum_{i=1}^k\sigma_iu_iv_i^\top = U_k\Sigma_kV_k^\top$.

\begin{theorem}[Eckart--Young--Mirsky]\label{svd:thm:eckart-young}
Among all matrices $B$ of rank at most $k$,
\[
  \norm{A - A_k}_F = \min_{\mathrm{rank}(B)\leq k}\norm{A - B}_F = \Bigl(\sum_{i>k}\sigma_i^2\Bigr)^{1/2},
  \qquad
  \norm{A - A_k}_2 = \min_{\mathrm{rank}(B)\leq k}\norm{A - B}_2 = \sigma_{k+1}.
\]
\end{theorem}

\begin{proof}[Proof for the spectral norm]
That $\norm{A - A_k}_2 = \sigma_{k+1}$ follows from $A - A_k = \sum_{i>k}\sigma_iu_iv_i^\top$.
Let $B$ have rank $\leq k$; then $\ker(B)$ has dimension $\geq n-k$ and meets the
$(k+1)$-dimensional space $\operatorname{span}(v_1,\dots,v_{k+1})$ in some unit vector $x$.
Then $\norm{(A-B)x} = \norm{Ax} = \bigl(\sum_{i\leq k+1}\sigma_i^2(v_i^\top x)^2\bigr)^{1/2}\geq\sigma_{k+1}$.
The Frobenius case is proved in \cref{ch:pca-overview} in the language of PCA.
\end{proof}

\begin{example}[Rank-one approximation]
In \cref{svd:ex:2x2},
$A_1 = \sigma_1u_1v_1^\top = 3\sqrt5\cdot\frac1{\sqrt{10}}\begin{psmallmatrix}1\\3\end{psmallmatrix}\frac1{\sqrt2}\begin{psmallmatrix}1&1\end{psmallmatrix} = \begin{psmallmatrix}1.5&1.5\\4.5&4.5\end{psmallmatrix}$,
and $A - A_1 = \begin{psmallmatrix}1.5&-1.5\\-0.5&0.5\end{psmallmatrix}$ has
$\norm{A - A_1}_F^2 = 5 = \sigma_2^2$: the rank-one matrix keeps $45/50 = 90\%$ of the
``energy'' $\norm A_F^2$.
\end{example}

Low-rank approximation is the principle behind PCA (\cref{ch:pca-overview}), latent
semantic analysis, image compression and the matrix factorization models of
recommender systems (\cref{ch:matrix-factorization}).

\section{The pseudoinverse}\label{svd:sec:pseudoinverse}

\begin{definition}[Moore--Penrose pseudoinverse]
With $A = U\Sigma V^\top$, the pseudoinverse is $A^+ = V\Sigma^+U^\top$, where
$\Sigma^+\in\R^{n\times m}$ is obtained by transposing $\Sigma$ and replacing each non-zero
$\sigma_i$ by $1/\sigma_i$: $A^+ = \sum_{i=1}^r\sigma_i^{-1}\,v_iu_i^\top$.
\end{definition}

Since $U^\top u_i = e_i$, the pseudoinverse undoes $A$ on the singular directions:
\[
  Av_i = \sigma_iu_i\qquad\text{and}\qquad A^+u_i = \frac1{\sigma_i}v_i\quad(i\leq r),
  \qquad A^+u_i = 0\quad(i>r).
\]
It is a true inverse between $\operatorname{span}(v_1,\dots,v_r) = \ker(A)^\perp$ and
$\operatorname{span}(u_1,\dots,u_r) = \operatorname{Im}(A)$, and sends
$\operatorname{Im}(A)^\perp$ to $0$. When $A$ is square and invertible, $A^+ = A^{-1}$; when
its columns are independent, $A^+ = (A^\top A)^{-1}A^\top$.

\begin{theorem}[Minimum-norm least squares]\label{svd:thm:minnorm}
$\hat x = A^+y$ solves the normal equations $A^\top A\hat x = A^\top y$, and has the smallest
norm among all least-squares solutions. Moreover $AA^+ = U_rU_r^\top$ is the orthogonal
projection onto $\operatorname{Im}(A)$.
\end{theorem}

\begin{proof}
$A^\top A A^+ = V\Sigma^\top U^\top U\Sigma V^\top V\Sigma^+U^\top = V(\Sigma^\top\Sigma\Sigma^+)U^\top = V\Sigma^\top U^\top = A^\top$,
because $\Sigma^\top\Sigma\Sigma^+ = \Sigma^\top$ (check on the diagonal: $\sigma_i^2\cdot\sigma_i^{-1} = \sigma_i$, and
$0$ where $\sigma_i = 0$). So $A^\top A(A^+y) = A^\top y$. All solutions are $A^+y + z$ with
$z\in\ker(A)$, and $A^+y\in\operatorname{span}(v_1,\dots,v_r) = \ker(A)^\perp$, so
$\norm{A^+y + z}^2 = \norm{A^+y}^2 + \norm z^2$ is smallest for $z = 0$. Finally
$AA^+ = U\Sigma\Sigma^+U^\top = \sum_{i\leq r}u_iu_i^\top$.
\end{proof}

\begin{example}[A rank-deficient system]
$A = \begin{psmallmatrix}1&1\\1&1\end{psmallmatrix}$ has rank one: $\sigma_1 = 2$, $u_1 = v_1 = \frac1{\sqrt2}(1,1)$.
Hence $A^+ = \frac12v_1u_1^\top = \frac14\begin{psmallmatrix}1&1\\1&1\end{psmallmatrix}$. For $y = (1,0)$,
every $x$ with $x_1 + x_2 = \frac12$ is a least-squares solution (the projection of $y$ on
$\operatorname{Im}(A)$ is $(\frac12,\frac12)$); $A^+y = (\frac14,\frac14)$ is the one of minimum norm.
\end{example}

\begin{chaptersummary}
\begin{itemize}[nosep]
  \item $A = U\Sigma V^\top = \sum_i\sigma_iu_iv_i^\top$ with orthogonal $U$, $V$ and singular values
        $\sigma_1\geq\dots\geq\sigma_r>0$; $v_i$, $\sigma_i^2$ are eigenpairs of $A^\top A$ and
        $u_i = Av_i/\sigma_i$.
  \item Geometrically: rotate ($V^\top$), stretch and change dimension ($\Sigma$), rotate ($U$);
        the unit sphere becomes an ellipsoid with semi-axes $\sigma_iu_i$, and $\norm A_2 = \sigma_1$.
  \item Truncating the sum gives the best rank-$k$ approximation, with errors
        $(\sum_{i>k}\sigma_i^2)^{1/2}$ in Frobenius norm and $\sigma_{k+1}$ in spectral norm.
  \item $A^+ = \sum_i\sigma_i^{-1}v_iu_i^\top$ gives the minimum-norm least-squares solution, and
        $AA^+$ is the orthogonal projection onto $\operatorname{Im}(A)$.
\end{itemize}
\end{chaptersummary}

\section*{Exercises}
\addcontentsline{toc}{section}{Exercises}

\begin{exercise}\label{svd:exr:1}
Show that $\norm A_F^2 = \sum_i\sigma_i^2$ and $\abs{\det A} = \prod_i\sigma_i$ for square $A$.
\solutionlink{svd:exr:1}
\end{exercise}

\begin{exercise}\label{svd:exr:2}
Compute the SVD of the $3\times2$ matrix $A = \begin{psmallmatrix}1&1\\1&1\\0&0\end{psmallmatrix}$ and its
pseudoinverse.
\solutionlink{svd:exr:2}
\end{exercise}

\begin{exercise}\label{svd:exr:3}
Verify the four Penrose conditions for $A^+$: $AA^+A = A$, $A^+AA^+ = A^+$, and $AA^+$, $A^+A$
symmetric.
\solutionlink{svd:exr:3}
\end{exercise}
\end{notechapter}
